
ce qui
montre que pour @@M0@@, @@M1@@ s'envoie sur la @@M2@@-ième fonction symétrique des @@M3@@, et par suite que @@M4@@ est injectif et identifie @@M5@@ à l'anneau des polynômes symétriques en les @@M6@@.

On a alors @@M7@@, et @@M8@@. Les @@M9@@ sont des polynômes par rapport aux @@M10@@~; si pour certain indice @@M11@@ on fait @@M12@@, on voit que @@M13@@ devient @@M14@@, donc que les @@M15@@ s'annulent pour @@M16@@, donc qu'ils contiennent @@M17@@ en facteur. Donc ils contiennent le produit des @@M18@@ en facteur, i.e. ils sont divisibles par @@M19@@ dans @@M20@@, donc aussi dans @@M21@@, le quotient de deux polynômes symétriques étant évidemment symétrique.

5.3. Comme @@M22@@ est intègre, on définit sans ambiguïté un élément @@M23@@ en posant (@@M24@@) :
@@M25@@

L'élément @@M26@@ est donc un polynôme par rapport aux @@M27@@, @@M28@@, isobare de poids @@M29@@ par rapport aux @@M30@@, car il en est de même de @@M31@@ d'après 2.3.

Soient @@M32@@ un entier @@M33@@, et @@M34@@ la @@M35@@-@@M36@@-algèbre engendrée par les générateurs @@M37@@ et @@M38@@ soumis aux relations @@M39@@ pour @@M40@@~; d'après la proposition 4.7, il existe un @@M41@@-homomorphisme et un seul @@M42@@, envoyant @@M43@@ sur @@M44@@ et @@M45@@ sur @@M46@@. De la définition (5.3.1) résulte que @@M47@@ pour tout @@M48@@.

Si @@M49@@ est un @@M50@@-anneau, et @@M51@@ un élément de @@M52@@ tel que @@M53@@ pour @@M54@@, on peut, pour tout @@M55@@, définir un élément @@M56@@ par substitution, ou encore en utilisant l'unique @@M57@@-homomorphisme de @@M58@@ dans 
